\section{Introduction}
\label{sec:intro}

Variational quantum eigensolvers have become a cornerstone of
near-term quantum simulation, trading circuit depth for a classical
optimization loop. For topological phases, however, the standard
variational workflow faces an extra hurdle: the order parameters that
theorists prefer are often the hardest to measure on hardware.

Topological phases on real hardware have been pursued both by direct
construction and by variational optimization. On superconducting
processors, Satzinger~\textit{et al.}~\cite{satzinger2021realizing}
prepared toric-code ground states and measured topological entanglement
entropy, demonstrating topological order itself as the target; direct
constructions of this kind yield satisfactory fidelities, however, they
are only applicable to exactly solvable fixed-point states.
For variational approaches, Sun~\textit{et al.}~\cite{sun2023efficient}
designed initialization-plus-variational ansatzes for symmetry-protected
topological phases and verified string order and edge modes on IBM
devices, although the ansatz structure has largely been determined by
trial and error for each target phase.

While the main idea of phase-matched initialization is rather simple,
its full potential for generic non-integrable models---with a single
ansatz family covering competing orders and with markers that survive
the transition to hardware---has not been uncovered.
The neighboring literatures we draw on are
established: one-dimensional symmetry-protected phases are classified by
projective representations and matrix-product
states~\cite{chen2011classification,schuch2011classifying};
randomized measurements reconstruct entanglement and many-body
invariants from few
settings~\cite{huang2020predicting,elben2022toolbox};
variational algorithms and their trainability limits are
surveyed~\cite{cerezo2021variational,tilly2022variational,mcclean2018barren};
error mitigation spans probabilistic cancellation and zero-noise
extrapolation~\cite{temme2017error,kandala2019error}; and variational
hardware experiments range from molecular spectra to spin
chains~\cite{colless2018robust,yu2023simulating}.
The two diagnostics we use have dedicated literatures of their own.
The dynamical spin structure factor of the anisotropic spin-$1/2$
Heisenberg chain was computed by Pereira~\textit{et
al.}~\cite{pereira2006dynamical}, establishing the momentum-resolved
response whose static limit underlies our $S(\pi)$ AFM marker.
String order was introduced for valence-bond
phases~\cite{dennijs1989preroughening}, characterized symmetry-wise for
quantum spin lattices~\cite{perezgarcia2008string}, and observed via
correlated particle-hole pairs in low-dimensional Mott
insulators~\cite{endres2011observation}; our string operator is the
directly measurable counterpart tracked across the same sweep.
On hardware, limited connectivity makes qubit mapping a first-class
problem~\cite{li2019tackling}; our stabilizer-plus-readout screening
answers the same question measurement-first, ranking the subgraphs the
experiment will actually use. Readout mitigation spans correlated
unfolding~\cite{bravyi2021mitigating,nation2021scalable} and model-free
expectation-value correction~\cite{vandenberg2022modelfree}; our
per-qubit $2\times2$ unfolding with nonnegativity sits in this
tensor-product family by construction. Finite-size gap scaling is
standard practice in spin-chain
numerics~\cite{sandvik2010computational}, with the Haldane-gap problem
as its classic case~\cite{haldane1983nonlinear,affleck1987rigorous,affleck1989quantum};
our $L=8$ assignments with gap-scaling checks follow that discipline,
with full curves in the Appendix.
In this work, we extend phase-matched variational preparation to the
dimerized XXZ chain, where trivial, topological, and antiferromagnetic
phases compete, and we characterize the prepared states with directly
measurable markers ($S(\pi)$, string) alongside the
reflection invariant $\tilde{Z}_{\mathcal R}$.
Our approach extends phase-matched preparation to a setting where three
orders compete within one model, while accounting for hardware
constraints from the start: premium-qubit screening, readout mitigation,
and markers restricted to directly measurable correlators.
Whether the prepared states track the simulator trends on hardware is
decided by the comparison in Sec.~\ref{sec:comparison}.

The paper is organized as follows. We first introduce the model and
its symmetries (Sec.~\ref{sec:model}) and map the finite-size phase
diagram (Sec.~\ref{sec:phasediagram}), with gap scaling in the Appendix.
We then define the directly measurable phase markers
(Sec.~\ref{sec:markers}) and prepare the variational states
(Sec.~\ref{sec:ansatz}), bring the protocol to hardware
(Sec.~\ref{sec:hardware}), and compare hardware with simulator results
(Sec.~\ref{sec:comparison}).
